% !TeX root = ../main.tex
\section{Recursive types}

Try to write down the type of a linked list of integers using only what we've built so far ($\tyint$, $\tybool$, $\to$, products, and sums). You'll want something like ``either empty, or an integer paired with... another list''. But ``another list'' is the very thing we're trying to define — the type refers to itself. None of the type formers from \S3, \S5, or \S6 let you write an equation like
\[
  \alpha = \unit + (\tyint \times \alpha)
\]
and get back a type $\alpha$ satisfying it, for the same reason STLC couldn't type the Y combinator (\S3): our grammar of types only builds \emph{finite} trees, and this equation, expanded out, has no finite solution — $\alpha = \unit + \tyint\times(\unit+\tyint\times(\unit+\dots))$, forever.

\subsection{$\mu$ types: naming the fixed point}

The fix is to add a type former that names this self-reference directly instead of trying to expand it away: $\mu\alpha.\,\tau$, read ``the type $\alpha$ satisfying $\alpha = \tau$'', i.e., the (least) fixed point of the equation $\alpha = \tau[\alpha]$ where $\tau$ may itself mention $\alpha$.
\[
  \tau ::= \dots \mid \mu\alpha.\,\tau \mid \alpha
\]
With this, the list type from above is just
\[
  \mathbf{List}\,\tyint \;=\; \mu\alpha.\; \unit + (\tyint \times \alpha)
\]
read as: a list of integers is either $\unit$ (representing $\nil$) or a pair of an integer and \emph{another list} — and that recursive reference is exactly what $\alpha$, bound by the enclosing $\mu$, stands for.

The same trick handles any self-referential data. A binary tree of integers:
\[
  \mathbf{Tree} \;=\; \mu\alpha.\; \tyint + (\alpha \times \tyint \times \alpha)
\]
either a leaf holding an integer, or an internal node with a left subtree, a value, and a right subtree — again, "subtree" is just $\alpha$, tied back to the whole type by $\mu$.

\subsection{fold and unfold}

Having a type $\mu\alpha.\,\tau$ doesn't automatically tell you how to construct or destruct its values — for that we need to relate $\mu\alpha.\,\tau$ to its one-step \textbf{unfolding} $\tau[\mu\alpha.\,\tau / \alpha]$ (substitute the whole recursive type back in for $\alpha$ — exactly one layer of the infinite expansion above). In the \textbf{iso-recursive} presentation (the one virtually every real typed language uses), $\mu\alpha.\,\tau$ and its unfolding are treated as two \emph{different, isomorphic} types, connected by an explicit pair of coercions:
\[
\infer[\textsc{T-Fold}]
  {\Gamma \vdash \mathbf{fold}_{\mu\alpha.\tau}(e) : \mu\alpha.\,\tau}
  {\Gamma \vdash e : \tau[\mu\alpha.\tau / \alpha]}
\qquad
\infer[\textsc{T-Unfold}]
  {\Gamma \vdash \mathbf{unfold}(e) : \tau[\mu\alpha.\tau / \alpha]}
  {\Gamma \vdash e : \mu\alpha.\,\tau}
\]
$\mathbf{fold}$ takes a value shaped like one unrolling of the recursive type and packages it up as the recursive type itself (this is what a list \emph{constructor} like $\mathit{cons}$ does under the hood); $\mathbf{unfold}$ takes a value of the recursive type and reveals one layer of its structure so you can pattern-match on it (this is what happens when you inspect a list to ask "is it empty or a cons?"). They're inverses of each other, and every construction or destruction of a recursive value has to explicitly pass through one of them.

The alternative design is \textbf{equi-recursive} types: instead of adding $\mathbf{fold}$/$\mathbf{unfold}$ as explicit terms, you simply declare $\mu\alpha.\,\tau$ and $\tau[\mu\alpha.\tau/\alpha]$ to be \emph{literally the same type}, and let the type checker silently unfold as many times as needed during type equality checking. This is more convenient on paper — no fold/unfold clutter in the term syntax — but it comes at a real implementation cost: type equality is no longer syntactic, it's checking equivalence of (potentially) infinite trees, which needs its own dedicated algorithm (typically bisimulation-style, comparing the two trees co-inductively). Iso-recursive types push that complexity into two constructs living explicitly in the term language, where a garden-variety syntax-directed type checker (like the one from \S3) can handle them with no special machinery; equi-recursive types keep the term language clean but push the complexity into type equality itself, which is much harder to reconcile with type \emph{inference} in particular — this is a large part of why ML-family languages, which lean so heavily on inference, are iso-recursive.

\subsection{algebraic data types are recursive types}

If any of this looks suspiciously like ordinary sum-of-products data declarations, that's because it is — algebraic data types in Haskell, OCaml, and Rust are recursive types, with $\mathbf{fold}$/$\mathbf{unfold}$ hidden behind constructors and pattern matching.
\begin{lstlisting}[language=Haskell]
data List a = Nil | Cons a (List a)

data Tree a = Leaf a | Node (Tree a) a (Tree a)
\end{lstlisting}
\lstinline|Cons| \emph{is} $\mathbf{fold}_{\mu\alpha.\, \unit + (a\times\alpha)}$ specialized to the \lstinline|Cons| branch: it takes an $a$ and a $\mathbf{List}\,a$ (one unrolling of the type) and produces a $\mathbf{List}\,a$. Pattern-matching on a list (\lstinline|case xs of Nil -> ...; Cons h t -> ...|) \emph{is} $\mathbf{unfold}$ followed by a case split on the resulting sum type. The same correspondence holds in OCaml (\lstinline|type 'a list = Nil | Cons of 'a * 'a list|) and Rust (\lstinline|enum List<T> { Nil, Cons(T, Box<List<T>>) }| — note Rust needs the explicit \lstinline|Box| because, unlike Haskell and OCaml, it doesn't box recursive fields automatically, and a type can't have infinite size unboxed). Every time you write a recursive \lstinline|data|/\lstinline|enum| declaration, you are writing $\mu\alpha.\,\tau$ with sugar on top; the language's compiler is generating fold and unfold for you (as the automatically-derived constructors and pattern matches) instead of making you write them by hand.

\subsection{closing the loop: the occurs check, revisited}

Back in \S4, Algorithm~W's unification step included an \textbf{occurs check}: when unifying $\alpha$ with some type $\tau$, if $\alpha$ occurs free in $\tau$ itself, unification fails rather than producing an "infinite" substitution. We saw exactly this in \S4's worked example: trying to type $\lambda f.\,f\,f$ forces $\alpha = \alpha\to\beta$, and the occurs check rejects it, because $\alpha$ appears on both sides.

It's worth being precise about \emph{why} that check exists: it exists specifically to keep ordinary Hindley-Milner types finite. HM's grammar of types (like STLC's) only builds finite trees, so equations like $\alpha = \alpha\to\beta$ genuinely have no solution within that grammar, and the occurs check is just unification correctly reporting "no solution exists" rather than silently looping forever trying to construct one.

Recursive types don't contradict that check or work around it — they \emph{change the grammar} so that equations of that shape suddenly do have a solution, on purpose, and mark that solution with an explicit $\mu$. $\mu\alpha.\,\tau$ is precisely a named, well-founded stand-in for what would otherwise be a rejected, infinitely-expanding unification solution: instead of asking "what infinite type solves $\alpha = \tau$", $\mu\alpha.\,\tau$ says "here is a single, finite type \emph{whose meaning} is the solution to that equation" and it comes with fold/unfold as the only sanctioned way to move between a $\mu$-type and one layer of its unrolling. So recursive types are a deliberate, controlled escape hatch — introduced by an explicit type former and mediated by explicit coercions — not a loophole that quietly reopens the door the occurs check was built to close. Ordinary HM inference (lists, trees, and all) still runs the occurs check on every other type variable; it's only at the places you've explicitly written or derived a $\mu$ that self-reference becomes legal.

\subsection{Exercises}
\begin{itemize}
  \item Write out $\mathbf{List}\,\tyint = \mu\alpha.\,\unit + (\tyint\times\alpha)$'s unfolding one level, and check that a two-element list $[1,2]$ can be built using two applications of $\mathbf{fold}$ (one for each \lstinline|Cons|, one final one for \lstinline|Nil|).
  \item Give a $\mu$-type for a \emph{non-empty} list of booleans (i.e., no $\nil$ case — every list has at least one element). Hint: it needs a base case that isn't $\unit$.
  \item Using the binary tree type $\mathbf{Tree} = \mu\alpha.\,\tyint + (\alpha\times\tyint\times\alpha)$ from this section, write the typing derivation for $\mathbf{fold}$ applied to a single-leaf tree holding the integer $5$.
  \item Explain in your own words why Rust's \lstinline|enum List<T> { Nil, Cons(T, List<T>) }| (without \lstinline|Box|) fails to compile, connecting your answer to the fact that $\mu\alpha.\,\tau$ types are usually represented as pointers at runtime.
  \item Iso-recursive types require explicit \lstinline|fold|/\lstinline|unfold| (or, in practice, constructors and pattern matches); equi-recursive types require none, but need a more complex type-equality check. Which design would you expect a language with full, unannotated Hindley-Milner inference to prefer, and why? Which would you expect a language that only does local type \emph{checking} against programmer-supplied annotations (no global inference) to tolerate more easily?
\end{itemize}
